\chapter{INTRODUCCIÓN}
\label{cap:introduccion}

En la actualidad, la práctica del ejercicio físico en gimnasios ha cobrado gran
relevancia como parte de un estilo de vida saludable. Sin embargo, diversos
estudios han evidenciado que una gran parte de los usuarios abandonan sus
rutinas en los primeros meses, principalmente por la falta de orientación,
motivación y seguimiento personalizado (Gjestvang, 2020). Esta problemática
afecta especialmente a los principiantes, quienes enfrentan barreras como el
desconocimiento del uso correcto de las máquinas, el alto coste de entrenadores
personales y la ausencia de planificación estructurada.

En respuesta a esta necesidad, se han desarrollado diversas soluciones
tecnológicas que buscan mejorar la experiencia del usuario en el gimnasio, como
aplicaciones móviles con realidad aumentada (Cruz, 2024) y sistemas de
información para rutinas y dietas (Fernández, 2018). No obstante, muchas de
estas herramientas carecen de personalización profunda y seguimiento
inteligente.

Este proyecto propone el desarrollo de una aplicación móvil basada en realidad
aumentada e inteligencia artificial, orientada a guiar a los usuarios en el uso
correcto de máquinas de ejercicio, planificar rutinas personalizadas y realizar
un seguimiento del progreso físico. La solución busca democratizar el acceso a
la orientación profesional, reducir el abandono temprano y fomentar la
constancia en el entrenamiento físico, utilizando tecnologías modernas como
React Native, Firebase y TensorFlow Lite.

Con esto, se pretende contribuir no solo al bienestar individual, sino también
a la optimización de recursos en gimnasios, ofreciendo una herramienta
accesible, interactiva y eficaz para la autogestión del entrenamiento.

% =====================================================================
\section{Antecedentes}
\label{sec:antecedentes}

(Yépez, 2017) en su proyecto de grado cuyo objetivo principal es la creación de
un gimnasio que combine asesoramiento presencial y virtual, orientado a la nueva
disciplina llamada Street Workout, que consiste en utilizar el propio peso
corporal para ejercitarse.

(Gjestvang, 2020) se realizó un estudio longitudinal en Noruega con el objetivo
de identificar las principales barreras en la iniciación y adherencia al
ejercicio en gimnasios. El proyecto evidenció que entre el 40\,\% y el 65\,\% de
los nuevos usuarios abandonan la asistencia regular en los primeros meses,
principalmente por la falta de motivación, orientación adecuada y seguimiento
personalizado. Estos resultados destacan la importancia de diseñar
herramientas tecnológicas que fortalezcan la adherencia a programas de
ejercicio estructurados.

(Cruz, 2024) desarrolló un proyecto de grado desarrollado en la Universidad
Mayor de San Simón, una aplicación móvil para gimnasios que utiliza escaneo de
código QR con realidad aumentada para mejorar la experiencia de los usuarios y
la gestión de entrenamientos, que permite a los usuarios obtener instrucciones
detalladas en tiempo real sobre los ejercicios.

(Fernandez, 2018) en su proyecto de grado desarrolló un sistema de información
cuyo objetivo principal es brindar apoyo a los usuarios o clientes de los
gimnasios, proporcionándoles rutinas de ejercicios y dietas, para que sirva de
guía sobre cómo usar las máquinas y como administrar sus alimentos para obtener
el objetivo deseado.

(BodBot IA Entrenador Personal, 2016) es una plataforma de fitness avanzada
creada por Edward Laux, Jason Storm y Sergio Prado, que ofrece entrenamientos de
IA personalizados según los objetivos, equipos, capacidades físicas y
dificultad del usuario, esta plataforma proporciona una guía paso a paso para
cada ejercicio, contiene personalización de peso corporal y registro de
alimentación diario.

(Ejercicios en el Gimnasio, 2021) desarrollada por Leap Fitness Group, es una
plataforma fitness para principiantes y asiduos al gimnasio, que funciona como
un entrenador personalizado, creando rutinas, mostrando la postura correcta para
realizar ejercicios esto por medio de instrucciones visuales y literales
explícitas, siguiendo el progreso con estadísticas y gráficos claros.

% =====================================================================
\section{Análisis del problema}
\label{sec:problema}

\subsection{Descripción del problema}

Este proyecto surge para resolver las principales barreras que los usuarios de
gimnasios enfrentan, particularmente la falta de orientación adecuada y la falta
de estructura en sus entrenamientos. (Gjestvang, 2020) Estudios demuestran que
entre el 40\,\% y el 65\,\% de los nuevos usuarios abandonan el gimnasio en los
primeros meses, principalmente por la falta de guía y seguimiento
personalizado. Los principales problemas identificados incluyen:

\begin{itemize}
    \item \textbf{Falta de conocimiento sobre el uso adecuado de las máquinas y
    la postura correcta:} Esto lleva a errores que pueden causar lesiones y
    dificultar el progreso hacia los objetivos físicos.
    \item \textbf{Costo elevado de los entrenadores personales:} Muchas personas
    no pueden acceder a una guía profesional debido al costo asociado, lo cual
    deja a la mayoría sin un apoyo adecuado.
    \item \textbf{Falta de seguimiento del progreso:} La falta de un sistema para
    registrar y analizar el progreso personal desmotiva a los usuarios, ya que
    no tienen una forma clara de visualizar sus avances.
    \item \textbf{Carencia de planificación estructurada del entrenamiento:} Sin
    una rutina clara y bien planificada, muchos usuarios se sienten perdidos y no
    logran aprovechar el tiempo de entrenamiento, lo cual fomenta la
    desmotivación y el abandono.
\end{itemize}

\subsection{Definición del problema}

¿Cómo desarrollar una aplicación para gimnasios que ofrezca seguimiento rutinas
de ejercicio personalizadas con Inteligencia artificial?

% =====================================================================
\section{Objetivos}
\label{sec:objetivos}

\subsection{Objetivo general}

Desarrollar un sistema inteligente para el entrenamiento en gimnasios,
proporcionando una guía estructurada y seguimiento personalizado de los
ejercicios.

\subsection{Objetivos específicos}

\begin{itemize}
    \item Guiar sobre el uso correcto de máquinas.
    \item Instruir en la ejecución correcta de ejercicios específicos.
    \item Diseñar rutinas de entrenamiento personalizadas y estructuradas,
    adaptadas a los objetivos de cada usuario.
    \item Construir herramientas de planificación que permitan a los usuarios
    organizar sus entrenamientos de manera eficiente.
    \item Registrar la ingesta de alimentos con el fin de controlar y evaluar las
    calorías consumidas por el usuario.
    \item Implementar un sistema de seguimiento del progreso físico.
    \item Generar reportes personalizados de progreso físico.
\end{itemize}

% =====================================================================
\section{Área de conocimiento}
\label{sec:area}

\textbf{Área:} Programación Móvil

\textbf{Subárea:} Sistema de Información

% =====================================================================
\section{Alcance}
\label{sec:alcance}

\begin{itemize}
    \item El sistema funcionará en sistemas operativos Android.
    \item Interfaz intuitiva para el registro de usuarios.
    \item Guía sobre el uso correcto de máquinas de gimnasio.
    \item Registrar el progreso físico (peso, estatura, edad, masa corporal).
    \item Creación de rutinas personalizadas según los objetivos del usuario.
    \item El registro de alimentación se limitará solo al conteo calórico, es
    decir una calculadora de IMC.
    \item El sistema no reemplaza la supervisión profesional de entrenadores o
    nutricionistas.
    \item No se contempla la integración con dispositivos externos (smartwatch,
    banda de frecuencia cardíaca, entre otros).
\end{itemize}

% =====================================================================
\section{Justificación}
\label{sec:justificacion}

\subsection{Justificación técnica}

El proyecto es técnicamente viable gracias a la combinación de tecnologías
modernas y flexibles. Se empleará React Native con Expo para garantizar una
experiencia fluida en dispositivos Android, permitiendo un desarrollo
multiplataforma con una sola base de código.

La base de datos y autenticación estarán gestionadas por Firebase, lo que
asegura escalabilidad, sincronización en tiempo real y seguridad en el manejo de
datos de usuarios.

El componente de inteligencia artificial se implementará para generar rutinas
personalizadas, basándose en parámetros como peso, altura, edad, nivel de
experiencia y objetivos físicos. Este módulo podrá entrenarse mediante modelos
preexistentes o redes neuronales ligeras optimizadas para dispositivos móviles.

Además, se integrará un sistema de visualización interactiva con imágenes o
animaciones que muestran la postura correcta en cada ejercicio, lo cual aumenta
la comprensión y reduce el riesgo de lesiones.

En conjunto, estas tecnologías garantizan un sistema robusto, escalable y con
potencial de expansión futura.

\subsection{Justificación económica}

La aplicación representa una alternativa de bajo costo frente a los servicios
tradicionales de entrenamiento personalizado. Al automatizar la generación de
rutinas mediante inteligencia artificial, se reducen los costos asociados a la
contratación de entrenadores físicos, permitiendo que los usuarios obtengan
orientación profesional sin incurrir en gastos adicionales.

Para los gimnasios, la herramienta puede convertirse en un complemento
estratégico que optimice el tiempo de los entrenadores y mejore la retención de
clientes, al ofrecerles una plataforma moderna de autogestión del
entrenamiento. Asimismo, el desarrollo de esta aplicación se apoya en
tecnologías de código abierto (React Native, Firebase, TensorFlow Lite, etc.),
lo que disminuye considerablemente los costos de desarrollo y mantenimiento del
sistema.

\subsection{Justificación social}

Este proyecto busca eliminar las barreras que limitan la práctica constante del
ejercicio físico, ofreciendo orientación y motivación a través de la tecnología.
Está dirigido principalmente a usuarios principiantes o personas que no tienen
acceso a un entrenador personal. Contribuye a la promoción de un estilo de vida
saludable, fomenta la constancia y previene lesiones al enseñar el uso correcto
de las máquinas.
